% AA-67
\chapter{Moduły Verma i klasyfikacja reprezentacji nieprzywiedlnych}
\label{aa:ch:lie-highest-existence}

W AA-65 wykazaliśmy, że waga najwyższa reprezentacji nieprzywiedlnej jest całkowita i dominująca oraz jednoznacznie wyznacza reprezentację z dokładnością do izomorfizmu. Pozostało wykazać istnienie. Zaczniemy od modułu, który może mieć nieskończony wymiar, a następnie nałożymy relacje zapewniające skończoność wymiaru.

W rozdziale $K$ jest algebraicznie domkniętym ciałem charakterystyki zero, $\mathfrak g$ niezerową skończenie wymiarową półprostą algebrą Liego nad $K$, $\mathfrak h$ jej maksymalną podalgebrą toralną, $\alpha_1,\ldots,\alpha_r$ bazą układu korzeni, a $e_i,f_i,h_i$ generatorami z AA-61. Stosujemy oznaczenia $P,P_+$ i $\preceq$ z Definicji~\ref{aa:def:lie-dominant-integral}. Dla modułów nieskończenie wymiarowych przestrzeń wagową definiujemy tym samym warunkiem $tv=\mu(t)v$ dla każdego $t\in\mathfrak h$.

\section{Moduły Verma}

\begin{twierdzenie}
\label{aa:tw:verma-universal}
Niech $K$ będzie algebraicznie domkniętym ciałem charakterystyki zero, $\mathfrak g$ niezerową skończenie wymiarową półprostą algebrą Liego nad $K$, $\mathfrak h$ maksymalną podalgebrą toralną, a $e_i,f_i,h_i$ generatorami z Twierdzenia~\ref{aa:tw:lie-simple-generators}, odpowiadającymi korzeniom prostym $\alpha_i$ ($1\le i\le r$). Dla każdego $\lambda\in\mathfrak h^*$ istnieje moduł $M(\lambda)$ generowany przez niezerowy wektor $v_\lambda$ taki, że $e_i v_\lambda=0$ i $h_i v_\lambda=\lambda(h_i)v_\lambda$. Każdy moduł $V$ generowany przez wektor $v$ spełniający te relacje jest obrazem dokładnie jednego homomorfizmu $M(\lambda)\to V$ wysyłającego $v_\lambda$ na $v$. Moduł $M(\lambda)$ jest sumą prostą przestrzeni wagowych o skończonych wymiarach, wszystkie jego wagi są $\preceq\lambda$, a $M(\lambda)_\lambda=Kv_\lambda$.
\end{twierdzenie}
\begin{proof}
Z Twierdzenia~\ref{aa:tw:serre-presentation} algebra $\mathfrak g$ ma prezentację Serre'a dla macierzy $A=(\alpha_j(h_i))$. Stosujemy konstrukcję z dowodu Twierdzenia~\ref{aa:tw:serre-highest-module} do $\lambda_i=\lambda(h_i)$. Otrzymujemy $M(\lambda)=T/J$, gdzie $T=K\langle x_1,\ldots,x_r\rangle$, $J$ jest dwustronnym ideałem generowanym przez ujemne relacje Serre'a $(\operatorname{ad}x_i)^{1-a_{ij}}x_j$ dla $i\ne j$, a $v_\lambda=1+J$. Ideał jest jednorodny względem wielostopnia $\mathbf n=(n_1,\ldots,n_r)$, określającego liczbę wystąpień poszczególnych liter. Nie zawiera jedynki. Składowa wielostopnia $\mathbf n$ jest przestrzenią wagową dla $\lambda-\sum_i n_i\alpha_i$, ponieważ tak działają $h_j$ w konstrukcji z AA-62. Różne wielostopnie dają różne wagi dzięki niezależności korzeni prostych i charakterystyce zero. Każda składowa ma skończony wymiar, gdyż istnieje tylko skończenie wiele słów danego wielostopnia. Stopień zerowy daje dokładnie $Kv_\lambda$.

Dla modułu $V$ z tezy definiujemy odwzorowanie $T\to V$, wysyłając słowo $x_{i_1}\cdots x_{i_k}$ na $f_{i_1}\cdots f_{i_k}v$, a jedynkę na $v$. Relacje Serre'a dla $f_i$ działają zerowo na $V$, więc odwzorowanie zeruje każdy element $uS_{ij}w$ generujący liniowo $J$. Indukcja po długości słowa, wykorzystująca $[h_i,f_j]=-a_{ij}f_j$ i $[e_i,f_j]=\delta_{ij}h_i$, pokazuje zgodność z działaniem $h_i$ oraz $e_i$: dla $e_i$ rekurencja jest dokładnie rekurencją z konstrukcji AA-62. Zgodność z $f_i$ wynika z mnożenia z lewej. Indukowane odwzorowanie $M(\lambda)\to V$ jest zatem homomorfizmem $\mathfrak g$-modułów. Jego obraz jest podmodułem zawierającym generator $v$, więc jest całym $V$. Jednoznaczność wynika z tego, że słowa w $f_i$ zastosowane do $v_\lambda$ rozpinają $M(\lambda)$.
\end{proof}

Moduł z tego twierdzenia nazywamy modułem Verma o wadze najwyższej $\lambda$.

\begin{lemat}
\label{aa:lem:verma-simple-quotient}
Niech $K$ będzie algebraicznie domkniętym ciałem charakterystyki zero, $\mathfrak g$ niezerową skończenie wymiarową półprostą algebrą Liego nad $K$, $\mathfrak h$ jej maksymalną podalgebrą toralną, $\alpha_1,\ldots,\alpha_r$ bazą korzeni, a $\lambda\in\mathfrak h^*$. Niech $M(\lambda)$ będzie modułem Verma o generatorze $v_\lambda$ z Twierdzenia~\ref{aa:tw:verma-universal}. Istnieje dokładnie jeden największy właściwy podmoduł $R(\lambda)\subset M(\lambda)$. Iloraz $L(\lambda)=M(\lambda)/R(\lambda)$ jest nieprzywiedlny, a obraz $v_\lambda$ jest niezerowym generatorem o wadze najwyższej $\lambda$.
\end{lemat}
\begin{proof}
Każdy podmoduł $N$ jest sumą prostą przecięć z przestrzeniami wagowymi. Istotnie, dla $u\in N$ rozkład $u=\sum_{\mu\in S}u_\mu$ ma skończony zbiór wag $S$. Wybieramy $t\in\mathfrak h$ takie, że wartości $\mu(t)$ dla $\mu\in S$ są parami różne. Jest to możliwe, gdyż skończenie wiele właściwych hiperpłaszczyzn nie pokrywa przestrzeni nad nieskończonym ciałem. Wielomian interpolacyjny
\[
 p_\mu(X)=\prod_{\nu\in S\setminus\{\mu\}}
 \frac{X-\nu(t)}{\mu(t)-\nu(t)}
\]
spełnia $p_\mu(t)u=u_\mu$, więc $u_\mu\in N$. Jeśli $N$ jest właściwy, jego składowa wagi $\lambda$ jest zerowa: w przeciwnym razie zawierałby generator $v_\lambda$. Suma wszystkich właściwych podmodułów, rozumiana jako zbiór skończonych sum ich elementów, ma nadal zerową składową wagi $\lambda$. Jest więc właściwym podmodułem $R(\lambda)$ zawierającym wszystkie właściwe podmoduły. Każdy właściwy podmoduł ilorazu ma przeciwobraz właściwy w $M(\lambda)$, zatem przeciwobraz jest zawarty w $R(\lambda)$ i podmoduł ilorazu jest zerowy. Generator nie znika w ilorazie.
\end{proof}

\section{Relacje zapewniające lokalną nilpotentność}

\begin{lemat}
\label{aa:lem:verma-integrable}
Niech $K$ będzie algebraicznie domkniętym ciałem charakterystyki zero, $\mathfrak g$ niezerową skończenie wymiarową półprostą algebrą Liego nad $K$, $\mathfrak h$ jej maksymalną podalgebrą toralną, a $e_i,f_i,h_i$ generatorami z Twierdzenia~\ref{aa:tw:lie-simple-generators}, odpowiadającymi korzeniom prostym $\alpha_i$ ($1\le i\le r$). Niech $\lambda\in P_+$, gdzie $P_+$ jest zbiorem wag całkowitych dominujących względem tej bazy korzeni. Oznaczmy $m_i=\lambda(h_i)\in\mathbb Z_{\ge0}$. W module Verma $M(\lambda)$ niech $N$ będzie sumą podmodułów generowanych przez $f_i^{m_i+1}v_\lambda$ dla $1\le i\le r$. Iloraz $Q=M(\lambda)/N$ jest niezerowy, ma jednowymiarową przestrzeń wagi $\lambda$, wszystkie jego wagi są $\preceq\lambda$, a każdy operator $e_i$ i $f_i$ działa na $Q$ lokalnie nilpotentnie (dla każdego wektora pewna potęga operatora zeruje ten wektor).
\end{lemat}
\begin{proof}
Wektor $f_i^{m_i+1}v_\lambda$ jest zerowany przez wszystkie $e_j$. Dla $j\ne i$ wynika to z $[e_j,f_i]=0$, a dla $j=i$ ze wzoru z Lematu~\ref{aa:lem:sl2-commutation}:
\[
 e_i f_i^{m_i+1}v_\lambda
 =(m_i+1)(m_i-(m_i+1)+1)f_i^{m_i}v_\lambda=0.
\]
Jeżeli wektor ten jest niezerowy, ma wagę $\lambda-(m_i+1)\alpha_i$ i generuje moduł będący obrazem modułu Verma tej wagi, na mocy Twierdzenia~\ref{aa:tw:verma-universal}. Wszystkie wagi tego podmodułu są więc $\preceq\lambda-(m_i+1)\alpha_i$; żadna nie jest równa $\lambda$. Jeśli wektor jest zerowy, generuje moduł zerowy. Wobec tego $N$ nie zawiera $v_\lambda$, a z rozkładu wagowego (argument projekcji z Lematu~\ref{aa:lem:verma-simple-quotient}) wynika, że $Q$ ma rozkład wagowy, zachowuje jednowymiarową składową $\lambda$ i ma skończenie wymiarowe przestrzenie wagowe.

Oznaczmy przez $v$ obraz $v_\lambda$ w $Q$. Operator $e_i$ podwyższa wagę o $\alpha_i$, więc na wektorze wagi $\lambda-\sum_j n_j\alpha_j$ jego potęga $n_i+1$ znika: odpowiadałaby ujemnemu współczynnikowi $n_i$ w różnicy od $\lambda$. Każdy wektor jest skończoną sumą wektorów wagowych, co daje lokalną nilpotentność.

Ustalmy $i$. Niech $\mathcal A$ będzie łączną podalgebrą $\operatorname{End}_K(Q)$ z jedynką, generowaną przez operatory $f_j$. Derywacja $D(a)=[f_i,a]$ zachowuje $\mathcal A$. Na generatorze $f_i$ znika, a na $f_j$ dla $j\ne i$ pewna jej potęga znika na mocy relacji Serre'a. Dla iloczynu $a_1\cdots a_k$ obowiązuje uogólniona reguła Leibniza
\[
 D^s(a_1\cdots a_k)
 =\sum_{s_1+\cdots+s_k=s}
 \frac{s!}{s_1!\cdots s_k!}D^{s_1}(a_1)\cdots D^{s_k}(a_k),
\]
wynikająca indukcyjnie z reguły Leibniza. Jeśli $D^{b_j}(a_j)=0$, to dla $s>\sum_j(b_j-1)$ każdy składnik znika. Stąd $D$ jest lokalnie nilpotentna na $\mathcal A$, również na skończonych kombinacjach słów.

Z równości $f_i a=D(a)+af_i$ otrzymujemy przez indukcję po $p$
\[
 f_i^p a v=\sum_{k=0}^{p}\binom pk D^k(a)f_i^{p-k}v.
\]
W kroku indukcyjnym stosujemy $f_iD^k(a)=D^{k+1}(a)+D^k(a)f_i$ i tożsamość Pascala. Dla $a\in\mathcal A$ wybierzmy $b$ z $D^b(a)=0$. Ponieważ $f_i^{m_i+1}v=0$, dla $p\ge b+m_i$ każdy składnik znika: albo $k\ge b$, albo $p-k\ge m_i+1$. Wektory $av$ rozpinają $Q$, więc $f_i$ działa lokalnie nilpotentnie.
\end{proof}

\section{Skończoność zbioru wag}

\begin{lemat}
\label{aa:lem:integrable-weight-symmetry}
Niech $K$ będzie algebraicznie domkniętym ciałem charakterystyki zero, $\mathfrak g$ niezerową skończenie wymiarową półprostą algebrą Liego nad $K$, $\mathfrak h$ jej maksymalną podalgebrą toralną, a $e_i,f_i,h_i$ generatorami z Twierdzenia~\ref{aa:tw:lie-simple-generators}, odpowiadającymi korzeniom prostym $\alpha_i$ ($1\le i\le r$). Niech $Q$ będzie $\mathfrak g$-modułem będącym algebraiczną sumą prostą przestrzeni wagowych, na którym $e_i,f_i$ działają lokalnie nilpotentnie dla każdego $i$. Zbiór wag $Q$ jest niezmienniczy względem odbić $s_i(\mu)=\mu-\mu(h_i)\alpha_i$.
\end{lemat}
\begin{proof}
Dla niezerowego wektora $u\in Q_\mu$ ustalmy $i$ i zapiszmy $e=e_i$, $f=f_i$, $h=h_i$. Rozważmy rozpiętość $S$ wektorów $f^a e^b u$ dla $a,b\ge0$. Tylko skończenie wiele $e^b u$ jest niezerowych, a dla każdego z nich tylko skończenie wiele $f^a e^b u$ jest niezerowych. Zatem $S$ jest skończenie wymiarowa. Każdy niezerowy wektor z tej rodziny jest wektorem wagowym dla $\mu+(b-a)\alpha_i$, więc $hS\subset S$. Niezmienniczość względem $f$ wynika wprost z definicji. Aby sprawdzić niezmienniczość względem $e$, korzystamy z tożsamości operatorowej
\[
 ef^a=f^a e+a f^{a-1}(h-a+1)\qquad(a\ge1).
\]
Dla $a=1$ jest to $ef=fe+h$. W kroku indukcyjnym mnożymy prawą stronę przez $f$ i używamy $(h-a+1)f=f(h-a-1)$, otrzymując $f^{a+1}e+(a+1)f^a(h-a)$. Zastosowanie wzoru do $e^b u$ daje kombinację $f^a e^{b+1}u$ i $f^{a-1}e^b u$. Dla $a=0$ niezmienniczość jest bezpośrednia.

Mamy więc skończenie wymiarowy moduł $\mathfrak{sl}_2$ rozpięty przez wektory, na których każdy $t\in\ker\alpha_i$ działa jako $\mu(t)$. Z Twierdzenia~\ref{aa:tw:sl2-complete-reducibility} wynika, że wraz z wartością własną $\mu(h_i)$ operatora $h$ występuje wartość przeciwna. Wybierzmy niezerowy wektor tej przeciwnej wartości w $S$. Ponieważ $\mathfrak h=Kh_i\oplus\ker\alpha_i$, jego waga jest dokładnie $s_i(\mu)$. Każde $s_i$ jest inwolucją, co daje niezmienniczość zbioru wag.
\end{proof}

\begin{twierdzenie}[Istnienie i klasyfikacja]
\label{aa:tw:lie-highest-classification}
Niech $K$ będzie algebraicznie domkniętym ciałem charakterystyki zero, $\mathfrak g$ niezerową skończenie wymiarową półprostą algebrą Liego nad $K$, $\mathfrak h$ maksymalną podalgebrą toralną, a $\alpha_1,\ldots,\alpha_r$ bazą korzeni. Dla każdej wagi całkowitej dominującej $\lambda\in P_+$ względem tej bazy korzeni moduł $L(\lambda)$ z Lematu~\ref{aa:lem:verma-simple-quotient} ma skończony wymiar. Przyporządkowanie $\lambda\mapsto[L(\lambda)]$ jest bijekcją między $P_+$ a klasami izomorfizmu niezerowych skończenie wymiarowych nieprzywiedlnych $\mathfrak g$-modułów.
\end{twierdzenie}
\begin{proof}
Niech $Q$ będzie niezerowym ilorazem z Lematu~\ref{aa:lem:verma-integrable}. Wystarczy wykazać, że $Q$ ma skończony wymiar, gdyż podmoduł $N$ jest właściwy, więc $N\subset R(\lambda)$ i $L(\lambda)$ jest ilorazem $Q$.

Niech $W$ będzie skończoną grupą Weyla układu korzeni, działającą na $\mathfrak h^*$ przez odbicia proste. Skończoność wynika z wiernego działania na skończonym zbiorze korzeni, które rozpinają $\mathfrak h^*$. Dla dowolnego $\mu\in\mathfrak h^*$ suma $\sum_{w\in W}w\mu$ jest niezmiennicza względem $W$ (mnożenie z lewej permutuje elementy grupy). Jeśli funkcjonał $\eta$ jest niezmienniczy, to $\eta=s_i\eta$ daje $\eta(h_i)\alpha_i=0$, czyli $\eta(h_i)=0$ dla każdego $i$. Ponieważ $h_i$ tworzą bazę $\mathfrak h$, mamy $\eta=0$. Zatem
\[
 \sum_{w\in W}w\mu=0.
\]
Jeśli $\mu$ jest wagą $Q$, to dla każdego $w\in W$ także $w\mu$ jest wagą na mocy Lematu~\ref{aa:lem:integrable-weight-symmetry}. Istnieją więc nieujemne liczby całkowite $n_j(w)$ takie, że
\[
 \lambda-w\mu=\sum_{j=1}^r n_j(w)\alpha_j.
\]
Sumując po $w$, dostajemy
\[
 |W|\lambda=\sum_{j=1}^r b_j\alpha_j,
 \qquad b_j=\sum_{w\in W}n_j(w).
\]
Współczynniki $b_j$ zależą wyłącznie od $\lambda$, nie od $\mu$, gdyż korzenie proste są liniowo niezależne. Są nieujemnymi liczbami całkowitymi (można obliczyć je już dla wagi $\mu=\lambda$). Dla elementu neutralnego $1\in W$ mamy $0\le n_j(1)\le b_j$. Zatem każda waga $\mu$ należy do skończonego zbioru
\[
 \left\{\lambda-\sum_j k_j\alpha_j:
 k_j\in\mathbb Z,\ 0\le k_j\le b_j\right\}.
\]
To ograniczenie dotyczy liczb całkowitych, nie porządku na $K$. Każda przestrzeń wagowa $Q$ ma skończony wymiar, więc $Q$ ma skończony wymiar.

Moduł $L(\lambda)$ jest niezerowy i nieprzywiedlny, a obraz $v_\lambda$ jest wektorem najwyższej wagi $\lambda$. Twierdzenie~\ref{aa:tw:lie-highest-necessary} mówi, że każda niezerowa skończenie wymiarowa reprezentacja nieprzywiedlna ma wagę najwyższą należącą do $P_+$. Twierdzenie~\ref{aa:tw:lie-highest-uniqueness} daje jej izomorfizm z $L(\lambda)$ oraz brak izomorfizmu dla różnych wag. Otrzymujemy żądaną bijekcję.
\end{proof}

\begin{wniosek}
\label{aa:cor:verma-integrable-simple}
Niech $K$ będzie algebraicznie domkniętym ciałem charakterystyki zero, $\mathfrak g$ niezerową skończenie wymiarową półprostą algebrą Liego nad $K$, $\mathfrak h$ jej maksymalną podalgebrą toralną, a $e_i,f_i,h_i$ generatorami odpowiadającymi korzeniom prostym $\alpha_i$ ($1\le i\le r$). Dla wagi całkowitej dominującej $\lambda$ oraz $m_i=\lambda(h_i)$ niech $N\subset M(\lambda)$ będzie sumą podmodułów generowanych przez $f_i^{m_i+1}v_\lambda$. Wówczas $M(\lambda)/N$ jest nieprzywiedlny, czyli $N=R(\lambda)$.
\end{wniosek}
\begin{proof}
W dowodzie Twierdzenia~\ref{aa:tw:lie-highest-classification} wykazaliśmy, że $Q=M(\lambda)/N$ ma skończony wymiar. Z Twierdzenia~\ref{aa:tw:lie-weyl-complete-reducibility} jest sumą prostą niezerowych modułów nieprzywiedlnych. Rzut obrazu $v_\lambda$ na każdy składnik jest niezerowy, gdyż w przeciwnym razie podmoduł generowany przez ten obraz byłby zawarty w sumie pozostałych składników. Każdy taki rzut ma wagę $\lambda$ i jest zerowany przez wszystkie $e_i$. Wobec tego każdy składnik wnosi niezerową składową do $Q_\lambda$. Lemat~\ref{aa:lem:verma-integrable} daje $\dim Q_\lambda=1$, więc występuje dokładnie jeden składnik. Moduł $Q$ jest nieprzywiedlny, a $N$ jest maksymalnym właściwym podmodułem $M(\lambda)$, zatem $N=R(\lambda)$.
\end{proof}

\begin{przyklad}[Rząd jeden]
\label{aa:prz:verma-sl2}
Dla $\mathfrak g=\mathfrak{sl}_2(K)$ oraz $\lambda(h)=m\in\mathbb Z_{\ge0}$ moduł Verma jest przestrzenią $K[x]$ z bazą $x^k$ dla $k\ge0$, gdyż nie ma ujemnych relacji Serre'a między różnymi generatorami. Z konstrukcji $fx^k=x^{k+1}$, $hx^k=(m-2k)x^k$, $ex^k=k(m-k+1)x^{k-1}$ dla $k\ge1$, a $e1=0$. Wektor $x^{m+1}$ jest zerowany przez $e$ i generuje podmoduł $x^{m+1}K[x]$. Iloraz ma bazę obrazów $1,x,\ldots,x^m$. Każdy niezerowy podmoduł ilorazu zawiera któryś z tych wektorów, dzięki jednowymiarowym różnym przestrzeniom wagowym. Kolejne zastosowania $e$ dają $1$, ponieważ współczynniki dla $1\le k\le m$ są niezerowe; następnie zastosowania $f$ dają całą bazę. Iloraz jest więc nieprzywiedlny i równy z dokładnością do izomorfizmu $L(m\omega_1)$, o wymiarze $m+1$. Odtworzyliśmy klasyfikację z AA-60.
\end{przyklad}

W połączeniu z Twierdzeniem~\ref{aa:tw:lie-weyl-complete-reducibility} otrzymujemy opis wszystkich skończenie wymiarowych reprezentacji jako sum prostych modułów $L(\lambda)$ dla $\lambda\in P_+$. Nie wyznaczaliśmy jeszcze ich wymiarów ani krotności poszczególnych wag w wyższych rzędach.
